\documentclass[a4paper,11pt]{article}

% Load ifthen package for conditional compilation
\usepackage{ifthen}

% Cover selection mechanism
% Options: none (default), 1, 2, etc. - numbers correspond to cover files in covers directory
\newcommand{\covertype}{none}

% ========== BRAND IDENTITY SETTINGS ==========
% Customize these to match your organization's branding
\newcommand{\companyName}{Lagen}
\newcommand{\companyTagline}{Professional Document Systems}
\newcommand{\documentVersion}{1.0.3}
\newcommand{\documentStatus}{Final Release}
\newcommand{\documentClassification}{Internal Use}

% Import all configuration and styling from config.tex
% ========== CORE CONFIGURATION ==========
% Lightweight core configuration for basic LaTeX documents
% ========== ESSENTIAL PACKAGES ==========
% Instructions: These packages handle font setup and character encoding
% Keep these as-is unless you need different fonts
\usepackage{inputenc}
\usepackage[T1]{fontenc}
\usepackage{lmodern}
\usepackage{ifthen}
\usepackage{geometry}
\usepackage{fontspec}
\usepackage{setspace}
\usepackage{xcolor}
\usepackage{morefloats}
\usepackage[table]{xcolor}
% Load TikZ for covers and templates
\usepackage{tikz}
\usepackage{eso-pic}
% Load fp package for mathematical calculations (needed by some covers)
\usepackage{fp}
% Load datetime2 package for date formatting (needed by covers 4 and 5)
\usepackage{datetime2}
\ifdefined\pdftexversion
  \usepackage{microtype}
\fi

% ========== BASIC PAGE LAYOUT ==========
\geometry{
    top=1.2in,
    bottom=1.2in,
    left=1.1in,
    right=1.1in,
    headheight=14pt,
    headsep=0.3in,
    footskip=0.45in
}

% ========== BASIC COLORS ==========
% Essential colors only - load brand_colors.tex if it exists
\definecolor{primary}{HTML}{0F172A}
\definecolor{secondary}{HTML}{334155}
\definecolor{accent}{HTML}{2563EB}
\definecolor{success}{HTML}{059669}
\definecolor{warning}{HTML}{D97706}
\definecolor{background}{HTML}{F8FAFC}
\definecolor{border}{HTML}{E2E8F0}
\definecolor{textgray}{HTML}{64748B}

% Try to load brand colors if available
\IfFileExists{brand_colors.tex}{
    % ========== BRAND COLORS ==========
% Define your brand colors here to maintain consistent branding across your document
% These colors will be used throughout the document for various elements

% Primary Brand Colors - CUSTOMIZE THESE
\definecolor{brandPrimary}{HTML}{2563EB}    % Primary brand color
\definecolor{brandSecondary}{HTML}{4F46E5}  % Secondary brand color
\definecolor{brandAccent}{HTML}{10B981}     % Accent color for highlights and emphasis

% Text Colors - Derived from brand colors
\definecolor{primary}{HTML}{0F172A}         % Main text color
\definecolor{secondary}{HTML}{334155}       % Secondary text color
\definecolor{accent}{HTML}{2563EB}          % Links and highlights (defaults to brandPrimary)

% UI Element Colors - These should work well with your brand colors
\definecolor{lightaccent}{HTML}{3B82F6}     % Lighter accent variant
\definecolor{success}{HTML}{059669}         % Success/positive elements
\definecolor{warning}{HTML}{D97706}         % Warning/attention elements
\definecolor{background}{HTML}{F8FAFC}      % Light background for boxes
\definecolor{border}{HTML}{E2E8F0}          % Borders and lines
\definecolor{textgray}{HTML}{64748B}        % Muted text
\definecolor{lightgray}{HTML}{F1F5F9}       % Table alternating rows

% Code Colors
\definecolor{codebg}{HTML}{F8F8F8}          % Code block background
\definecolor{codecomment}{HTML}{6A737D}     % Code comments
\definecolor{codestring}{HTML}{032F62}      % Code strings
\definecolor{codekey}{HTML}{D73A49}         % Code keywords

% Table Colors
\definecolor{cardbg}{HTML}{F7F9FC}          % Card background color
\definecolor{headerbg}{HTML}{4A90E2}        % Table header background (defaults to a shade of brandPrimary)
\definecolor{headertext}{HTML}{FFFFFF}      % Text color for table headers
\definecolor{rowalt}{HTML}{F1F4F9}          % Alternating row color
\definecolor{frameborder}{HTML}{D3DCE6}     % Border color for frames

% ========== COLOR CASCADE SYSTEM ==========
% This automatically adjusts certain colors based on your brand colors
% Uncomment and adjust these lines to automatically derive colors from your brand colors

% Table header background from brand primary
\colorlet{headerbg}{brandPrimary}

% Accent color from brand primary
\colorlet{accent}{brandPrimary}

% Light accent as a lighter shade of brand primary
\colorlet{lightaccent}{brandPrimary!75}

% Success color from brand accent
\colorlet{success}{brandAccent}

% Border color from brand colors
\colorlet{border}{brandPrimary!20}

% Frame border from brand secondary
\colorlet{frameborder}{brandSecondary!30}

% ========== END COLOR CASCADE SYSTEM ========== 
    \typeout{^^J INFO: Loaded custom brand colors from brand_colors.tex^^J}
}{
    \typeout{^^J INFO: Using basic colors (brand_colors.tex not found)^^J}
}

% ========== BASIC TYPOGRAPHY ==========
\setlength{\parindent}{0pt}
\setlength{\parskip}{8pt plus 2pt minus 1pt}

% Basic text commands
\newcommand{\brandname}[1]{\textsc{#1}}
\newcommand{\productname}[1]{\textbf{#1}}

% ========== CONDITIONAL LOADING SYSTEM ==========
% Define flags for optional features (set these before loading config-core.tex)
\providecommand{\EnableAdvancedTypography}{false}
\providecommand{\EnableMathematics}{false}
\providecommand{\EnableTables}{false}
\providecommand{\EnableBoxes}{false}
\providecommand{\EnableLists}{false}
\providecommand{\EnableCode}{false}
\providecommand{\EnableImages}{false}
\providecommand{\EnableAlgorithms}{false}
\providecommand{\EnableReferences}{false}

% Load additional modules based on flags
\ifthenelse{\equal{\EnableAdvancedTypography}{true}}{
    % ========== ADVANCED TYPOGRAPHY MODULE ==========
% Load this module only when advanced typography features are needed

% Advanced typography packages
\usepackage{fontspec}
\usepackage{setspace}
% Only load microtype for pdfTeX (not XeTeX/LuaTeX)
\ifdefined\pdftexversion
    \usepackage{microtype}
\fi
\usepackage{soul}
\usepackage{ragged2e}
\usepackage{textcase}
\usepackage{relsize}
\usepackage{scalefnt}
\usepackage{lettrine}
\usepackage[english]{babel}
\usepackage[autostyle=true,english=american]{csquotes}
\usepackage[all]{nowidow}
\usepackage{needspace}

% Font setup with Inter (if available)
\IfFontExistsTF{Inter}{
    \setmainfont{Inter}[
        UprightFont = *-Regular,
        BoldFont = *-SemiBold,
        ItalicFont = *-Italic,
        BoldItalicFont = *-SemiBoldItalic,
        LetterSpace = 3.0
    ]
    \setsansfont{Inter}[
        UprightFont = *-Regular,
        BoldFont = *-SemiBold,
        ItalicFont = *-Italic,
        BoldItalicFont = *-SemiBoldItalic,
        LetterSpace = 3.0
    ]
}{
    \typeout{^^J WARNING: Inter font not found, using default fonts^^J}
}

% Enhanced brand text macros
\renewcommand{\brandname}[1]{{\addfontfeatures{LetterSpace=12.0}\textsc{#1}}}
\renewcommand{\productname}[1]{{\addfontfeatures{LetterSpace=8.0}\textbf{#1}}}

% Line spacing commands
\newcommand{\setSpacingSingle}{\setstretch{1.0}}
\newcommand{\setSpacingOneHalf}{\setstretch{1.5}}
\newcommand{\setSpacingDouble}{\setstretch{2.0}}
\newcommand{\customspacing}[1]{\setstretch{#1}}

% Text decoration commands
\newcommand{\textunderline}[1]{\underline{#1}}
\newcommand{\textstrikeout}[1]{\st{#1}}
\newcommand{\texthi}[2][yellow]{\colorbox{#1}{#2}}

% Letter spacing adjustments
\newcommand{\wide}[1]{\so{#1}}
\newcommand{\wider}[1]{\sowide{#1}}
\newcommand{\narrow}[1]{\sonormal{#1}}
\newcommand{\letterspace}[2][100]{{\addfontfeatures{LetterSpace=#1}#2}}

% Font size adjustments
\newcommand{\textsm}[1]{\textsmaller{#1}}
\newcommand{\textlg}[1]{\textlarger{#1}}

% Case transformation
\newcommand{\uppercased}[1]{\MakeTextUppercase{#1}}
\newcommand{\lowercased}[1]{\MakeTextLowercase{#1}}

% Text alignment
\newcommand{\justifiedtext}{\justify}
\newcommand{\leftalignedtext}{\raggedright}
\newcommand{\rightalignedtext}{\raggedleft}
\newcommand{\centeredtext}{\centering}

% Drop cap configuration
\renewcommand{\LettrineFontHook}{\color{accent}\bfseries}
\setcounter{DefaultLines}{3}
\renewcommand{\DefaultLoversize}{0.1}
\renewcommand{\DefaultLraise}{0}

% Typography fine-tuning
\setlength{\emergencystretch}{3em}
\hyphenpenalty=500
\clubpenalty=10000
\widowpenalty=10000

% Microtype setup (only for pdfTeX)
\ifdefined\pdftexversion
    \microtypesetup{
        activate=true,
        final=true,
        factor=1100
    }
\fi

% Smart quotes setup
\MakeOuterQuote{"}

% Font feature control
\newcommand{\withfeatures}[2]{{\addfontfeatures{#1}#2}}

    \typeout{^^J INFO: Loaded advanced typography module^^J}
}{}

\ifthenelse{\equal{\EnableMathematics}{true}}{
    % ========== MATHEMATICS MODULE ==========
% Load this module only when mathematical typesetting is needed

% Core math packages
\usepackage{amsmath}
\usepackage{amssymb}
\usepackage{fixmath}
\usepackage{siunitx}
\usepackage{mathtools}
\usepackage{empheq}
\usepackage{cases}
\usepackage{amsthm}
\usepackage{stmaryrd}
\usepackage{textcomp}
\usepackage{physics}
\usepackage{tensor}
\usepackage{braket}
\usepackage{cancel}

% Configure siunitx
\sisetup{
    group-separator={,},
    group-minimum-digits=4,
    detect-weight=true,
    detect-family=true
}

% Define common units
\DeclareSIUnit\USD{\$}
\DeclareSIUnit\hour{h}
\DeclareSIUnit\year{yr}

% Theorem environments
\theoremstyle{definition}
\newtheorem{definition}{Definition}[section]
\newtheorem{theorem}{Theorem}[section]
\newtheorem{lemma}[theorem]{Lemma}
\newtheorem{corollary}[theorem]{Corollary}
\newtheorem{proposition}[theorem]{Proposition}
\newtheorem{example}{Example}[section]

\theoremstyle{remark}
\newtheorem{remark}{Remark}[section]
\newtheorem{note}{Note}[section]

% Math colors
\definecolor{mathblue}{RGB}{0,82,155}
\definecolor{mathred}{RGB}{204,0,0}
\definecolor{mathgreen}{RGB}{0,128,0}
\definecolor{mathpurple}{RGB}{128,0,128}

% Common mathematical sets
\newcommand{\R}{\mathbb{R}}
\newcommand{\C}{\mathbb{C}}
\newcommand{\N}{\mathbb{N}}
\newcommand{\Z}{\mathbb{Z}}
\newcommand{\Q}{\mathbb{Q}}
\newcommand{\F}{\mathbb{F}}

% Enhanced derivatives
\renewcommand{\d}{\mathrm{d}}
\newcommand{\dt}{\frac{\d}{\d t}}
\newcommand{\dx}{\frac{\d}{\d x}}
\newcommand{\dy}{\frac{\d}{\d y}}
\newcommand{\dz}{\frac{\d}{\d z}}
\newcommand{\prt}{\partial}
\newcommand{\pdx}[1]{\frac{\partial #1}{\partial x}}
\newcommand{\pdy}[1]{\frac{\partial #1}{\partial y}}
\newcommand{\pdz}[1]{\frac{\partial #1}{\partial z}}
\newcommand{\pdt}[1]{\frac{\partial #1}{\partial t}}

% Vector calculus
\providecommand{\grad}{\nabla}
\providecommand{\divg}{\nabla \cdot}
\providecommand{\curl}{\nabla \times}
\providecommand{\lapl}{\nabla^2}

% Vectors and matrices
\newcommand{\vect}[1]{\mathbf{#1}}
\newcommand{\mat}[1]{\mathbf{#1}}
\newcommand{\T}{\mathsf{T}}

% Probability and statistics
\newcommand{\prob}[1]{\mathrm{P}\left(#1\right)}
\newcommand{\E}[1]{\mathrm{E}\left[#1\right]}
\renewcommand{\var}[1]{\mathrm{Var}\left(#1\right)}
\newcommand{\cov}[1]{\mathrm{Cov}\left(#1\right)}
\newcommand{\normal}{\mathcal{N}}
\newcommand{\uniform}{\mathcal{U}}

% Math operators
\providecommand{\tr}{\operatorname{tr}}
\providecommand{\diag}{\operatorname{diag}}
\providecommand{\rank}{\operatorname{rank}}
\newcommand{\sign}{\operatorname{sign}}
\providecommand{\lcm}{\operatorname{lcm}}
\providecommand{\gcd}{\operatorname{gcd}}

% Boxed equations
\newcommand{\boxedeq}[1]{\begin{empheq}[box=\fbox]{equation}#1\end{empheq}}
\newcommand{\colorboxedeq}[2]{\begin{empheq}[box={\fcolorbox{#1}{white}}]{equation}#2\end{empheq}}

    \typeout{^^J INFO: Loaded mathematics module^^J}
}{}

\ifthenelse{\equal{\EnableTables}{true}}{
    % ========== TABLES MODULE ==========
% Load this module only when advanced table features are needed

% Table packages
\usepackage{booktabs}
\usepackage{tabularx}
\usepackage{caption}
\usepackage{longtable}
\usepackage{array}
\usepackage{multirow}

% Define enhanced column types
\newcolumntype{L}[1]{>{\raggedright\arraybackslash}p{#1}}
\newcolumntype{C}[1]{>{\centering\arraybackslash}p{#1}}
\newcolumntype{R}[1]{>{\raggedleft\arraybackslash}p{#1}}
\newcolumntype{N}{>{\raggedleft\arraybackslash}X}

% Table styling
\captionsetup{font=small, labelfont=bf, labelsep=colon}

% Additional table colors (if not already defined)
\providecommand{\definecolor}[3]{}
\definecolor{cardbg}{HTML}{F7F9FC}
\definecolor{headerbg}{HTML}{4A90E2}
\definecolor{headertext}{HTML}{FFFFFF}
\definecolor{rowalt}{HTML}{F1F4F9}
\definecolor{frameborder}{HTML}{D3DCE6}

    \typeout{^^J INFO: Loaded tables module^^J}
}{}

\ifthenelse{\equal{\EnableBoxes}{true}}{
    % ========== BOXES MODULE ==========
% Load this module only when callout boxes are needed

% Required packages
\usepackage{tikz}
\usetikzlibrary{shadows}
\usepackage[most]{tcolorbox}

% Info box style - use for general information
\newtcolorbox{infobox}[1][]{
    colback=background,
    colframe=accent,
    coltext=primary,
    boxrule=1pt,
    arc=3pt,
    left=12pt,
    right=12pt,
    top=8pt,
    bottom=8pt,
    enhanced,
    drop shadow,
    fonttitle=\bfseries\sffamily\color{white},
    title style={left color=accent, right color=accent!80!black},
    attach boxed title to top left={xshift=8pt, yshift=-2pt},
    #1
}

% Alert box style - use for cautions and important notes
\newtcolorbox{alertbox}[1][]{
    colback=warning!5,
    colframe=warning,
    coltext=primary,
    boxrule=1pt,
    arc=3pt,
    left=12pt,
    right=12pt,
    top=8pt,
    bottom=8pt,
    enhanced,
    drop shadow,
    fonttitle=\bfseries\sffamily\color{white},
    title style={left color=warning, right color=warning!80!black},
    attach boxed title to top left={xshift=8pt, yshift=-2pt},
    #1
}

% Success box style - use for positive highlights
\newtcolorbox{successbox}[1][]{
    colback=success!5,
    colframe=success,
    coltext=primary,
    boxrule=1pt,
    arc=3pt,
    left=12pt,
    right=12pt,
    top=8pt,
    bottom=8pt,
    enhanced,
    drop shadow,
    fonttitle=\bfseries\sffamily\color{white},
    title style={left color=success, right color=success!80!black},
    attach boxed title to top left={xshift=8pt, yshift=-2pt},
    #1
}

% Icon-enhanced boxes (optional - loads fontawesome and awesomebox)
\IfFileExists{fontawesome5.sty}{
    \usepackage{fontawesome5}
    \IfFileExists{awesomebox.sty}{
        \usepackage{awesomebox}
        \typeout{^^J INFO: Loaded icon-enhanced boxes (fontawesome5 + awesomebox)^^J}
    }{
        \typeout{^^J WARNING: awesomebox package not found, icon boxes unavailable^^J}
    }
}{
    \typeout{^^J WARNING: fontawesome5 package not found, icon boxes unavailable^^J}
}

% Logo-enhanced boxes (optional - loads bclogo)
\IfFileExists{bclogo.sty}{
    \usepackage{bclogo}
    \typeout{^^J INFO: Loaded logo-enhanced boxes (bclogo)^^J}
}{
    \typeout{^^J WARNING: bclogo package not found, logo boxes unavailable^^J}
}

    \typeout{^^J INFO: Loaded boxes module^^J}
}{}

\ifthenelse{\equal{\EnableLists}{true}}{
    % ========== LISTS MODULE ==========
% Load this module only when enhanced list features are needed

% List customization package
\usepackage{enumitem}

% Basic list styling
\setlist[itemize,1]{
    leftmargin=18pt,
    itemsep=3pt plus 1pt minus 1pt,
    parsep=0pt,
    topsep=6pt plus 2pt minus 2pt,
    label=\color{accent}$\bullet$
}
\setlist[itemize,2]{
    leftmargin=16pt,
    itemsep=2pt,
    parsep=0pt,
    topsep=3pt,
    label=\color{secondary}$\circ$
}
\setlist[itemize,3]{
    leftmargin=14pt,
    itemsep=1pt,
    parsep=0pt,
    topsep=2pt,
    label=\color{textgray}{\scriptsize$\blacksquare$}
}

% Enhanced enumerate styling
\setlist[enumerate,1]{
    leftmargin=22pt,
    itemsep=3pt plus 1pt minus 1pt,
    parsep=0pt,
    topsep=6pt plus 2pt minus 2pt,
    label=\color{accent}\arabic*.
}
\setlist[enumerate,2]{
    leftmargin=20pt,
    itemsep=2pt,
    parsep=0pt,
    topsep=3pt,
    label=\color{secondary}(\alph*)
}
\setlist[enumerate,3]{
    leftmargin=18pt,
    itemsep=1pt,
    parsep=0pt,
    topsep=2pt,
    label=\color{textgray}\roman*.
}

% Description list styling
\setlist[description]{
    font=\normalfont\color{accent}\bfseries,
    labelwidth=2cm,
    leftmargin=2.5cm,
    itemsep=4pt,
    parsep=0pt
}

% TikZ bullets for professional styling
\usepackage{tikz}
\newcommand{\tikzbullet}[1]{%
    \tikz[baseline=-0.7ex]\node[circle,fill=#1,inner sep=1.5pt]{};%
}

% Compact and expanded list variants
\newenvironment{compactlist}
    {\begin{itemize}[nosep]}
    {\end{itemize}}

\newenvironment{expandedlist}
    {\begin{itemize}[itemsep=8pt,parsep=4pt]}
    {\end{itemize}}

% Task lists with icons (requires fontawesome5)
\IfFileExists{fontawesome5.sty}{
    \usepackage{fontawesome5}
    
    \newlist{tasklist}{itemize}{1}
    \setlist[tasklist]{
        leftmargin=20pt,
        label=\color{success}\faCheck,
        itemsep=4pt,
        parsep=2pt
    }
    
    \newlist{pendingtask}{itemize}{1}
    \setlist[pendingtask]{
        leftmargin=20pt,
        label=\color{warning}\faHourglass,
        itemsep=4pt,
        parsep=2pt
    }
    
    \newlist{failedtask}{itemize}{1}
    \setlist[failedtask]{
        leftmargin=20pt,
        label=\color{accent}\faTimes,
        itemsep=4pt,
        parsep=2pt
    }
    
    % Priority list environments
    \newlist{highpriority}{itemize}{1}
    \setlist[highpriority]{
        leftmargin=20pt,
        label=\color{accent}\faExclamationCircle,
        itemsep=4pt
    }
    
    \newlist{mediumpriority}{itemize}{1}
    \setlist[mediumpriority]{
        leftmargin=20pt,
        label=\color{warning}\faExclamation,
        itemsep=4pt
    }
    
    \newlist{lowpriority}{itemize}{1}
    \setlist[lowpriority]{
        leftmargin=20pt,
        label=\color{success}\faInfoCircle,
        itemsep=4pt
    }
    
    \typeout{^^J INFO: Loaded task and priority lists (fontawesome5)^^J}
}{
    \typeout{^^J WARNING: fontawesome5 not found, task lists unavailable^^J}
}

    \typeout{^^J INFO: Loaded lists module^^J}
}{}

\ifthenelse{\equal{\EnableCode}{true}}{
    % ========== CODE MODULE ==========
% Load this module only when code listings are needed

% Code colors
\definecolor{codebg}{HTML}{F8F8F8}
\definecolor{codecomment}{HTML}{6A737D}
\definecolor{codestring}{HTML}{032F62}
\definecolor{codekey}{HTML}{D73A49}

% Basic code listings
\usepackage{listings}
\lstset{
    basicstyle=\ttfamily\small,
    backgroundcolor=\color{codebg},
    commentstyle=\color{codecomment},
    keywordstyle=\color{codekey}\bfseries,
    stringstyle=\color{codestring},
    numbers=left,
    numberstyle=\tiny\color{textgray},
    numbersep=10pt,
    tabsize=4,
    breaklines=true,
    breakatwhitespace=false,
    frame=single,
    rulecolor=\color{border},
    framesep=8pt,
    showstringspaces=false
}

% Advanced syntax highlighting with minted (optional)
\IfFileExists{minted.sty}{
    \usepackage{minted}
    \setminted{
        style=default,
        bgcolor=codebg,
        linenos=true,
        breaklines=true,
        tabsize=4,
        fontsize=\small,
        frame=single,
        framesep=8pt
    }
    \typeout{^^J INFO: Loaded advanced syntax highlighting (minted)^^J}
}{
    \typeout{^^J WARNING: minted package not found, using basic listings only^^J}
}

    \typeout{^^J INFO: Loaded code module^^J}
}{}

\ifthenelse{\equal{\EnableImages}{true}}{
    % ========== IMAGES MODULE ==========
% Load this module only when advanced image features are needed

% Image packages
\usepackage{graphicx}
\usepackage[export,graphics,noclipbox]{adjustbox}
\usepackage{varwidth}
\usepackage{tikz}

% Create fallback box for missing images
\newsavebox{\imagefallbackbox}

% Command for images with fallback text
\newcommand{\imageWithFallback}[4][]{%
  \begingroup
  \sbox{\imagefallbackbox}{%
    \begin{varwidth}{\textwidth}
      \centering
      \textcolor{textgray}{\textbf{#4}}
    \end{varwidth}
  }%
  \IfFileExists{#2}{%
    \includegraphics[#1,width=#3]{#2}%
  }{%
    \fcolorbox{border}{background}{%
      \begin{minipage}[c][0.75\ht\imagefallbackbox][c]{#3}
        \centering
        \textcolor{textgray}{\faImage}\\[0.3em]
        \usebox{\imagefallbackbox}
      \end{minipage}%
    }%
  }%
  \endgroup
}

% Rectangular image with border
\newcommand{\rectImage}[5][]{%
  \begin{figure}[htbp]
    \centering
    \fbox{\imageWithFallback[#1]{#2}{#3}{#5}}
    \ifx&#4&%
    \else
      \caption{#4}
    \fi
  \end{figure}
}

% Rounded corner image
\newcommand{\roundedImage}[6][]{%
  \begin{figure}[htbp]
    \centering
    \begin{tikzpicture}
      \node[inner sep=0pt] (image) {%
        \imageWithFallback[#1]{#2}{#3}{#6}%
      };
      \begin{scope}
        \clip[rounded corners=#4] (image.south west) rectangle (image.north east);
        \node[inner sep=0pt] {\imageWithFallback[#1]{#2}{#3}{#6}};
      \end{scope}
    \end{tikzpicture}
    \ifx&#5&%
    \else
      \caption{#5}
    \fi
  \end{figure}
}

% Circular image
\newcommand{\circularImage}[5][]{%
  \begin{figure}[htbp]
    \centering
    \begin{tikzpicture}
      \node[inner sep=0pt] (image) {%
        \imageWithFallback[#1]{#2}{#3}{#5}%
      };
      \begin{scope}
        \clip (image.center) circle (#3/2);
        \node[inner sep=0pt] {\imageWithFallback[#1]{#2}{#3}{#5}};
      \end{scope}
    \end{tikzpicture}
    \ifx&#4&%
    \else
      \caption{#4}
    \fi
  \end{figure}
}

% Image with drop shadow
\newcommand{\shadowImage}[5][]{%
  \begin{figure}[htbp]
    \centering
    \begin{tikzpicture}
      \node[inner sep=0pt, drop shadow={shadow xshift=2pt, shadow yshift=-2pt, opacity=0.5}] {%
        \imageWithFallback[#1]{#2}{#3}{#5}%
      };
    \end{tikzpicture}
    \ifx&#4&%
    \else
      \caption{#4}
    \fi
  \end{figure}
}

% Full-width responsive image
\newcommand{\fullwidthImage}[5][]{%
  \begin{figure}[htbp]
    \centering
    \makebox[\textwidth][c]{%
      \imageWithFallback[#1,height=#3,width=\textwidth,keepaspectratio]{#2}{\textwidth}{#5}%
    }
    \ifx&#4&%
    \else
      \caption{#4}
    \fi
  \end{figure}
}

    \typeout{^^J INFO: Loaded images module^^J}
}{}

\ifthenelse{\equal{\EnableAlgorithms}{true}}{
    % ========== ALGORITHMS MODULE ==========
% Load this module only when algorithm typesetting is needed

% Algorithm package
\usepackage[ruled,vlined]{algorithm2e}
\SetAlgoLined
\SetAlgoSkip{medskip}
\SetAlFnt{\footnotesize\sffamily}
\renewcommand{\algorithmcfname}{Algorithm}

    \typeout{^^J INFO: Loaded algorithms module^^J}
}{}

\ifthenelse{\equal{\EnableReferences}{true}}{
    % ========== REFERENCES MODULE ==========
% Load this module only when advanced references are needed

% Reference management packages
\usepackage[
    backend=biber,
    style=authoryear,
    sorting=nyt,
    maxbibnames=99,
    giveninits=true
]{biblatex}

% Hyperref and cross-references (order matters!)
\usepackage{hyperref}
\hypersetup{
    colorlinks=true,
    linkcolor=accent,
    urlcolor=accent,
    citecolor=accent,
    filecolor=accent,
    bookmarks=true,
    bookmarksopen=true,
    bookmarksnumbered=true,
    pdfstartview=FitH,
    pdfpagelayout=SinglePage
}

\usepackage{nameref}
\usepackage{varioref}
\usepackage{cleveref}

% Label formats
\labelformat{figure}{Figure~#1}
\labelformat{table}{Table~#1}
\labelformat{equation}{Equation~#1}
\labelformat{section}{Section~#1}

% Cleveref names
\crefname{figure}{Figure}{Figures}
\crefname{table}{Table}{Tables}
\crefname{equation}{Equation}{Equations}
\crefname{section}{Section}{Sections}

    \typeout{^^J INFO: Loaded references module^^J}
}{}

% ========== BASIC SECTION STYLING ==========
\usepackage{titlesec}
\titleformat{\section}
    {\color{primary}\fontsize{16}{20}\bfseries\sffamily}
    {\color{accent}\thesection}
    {1em}
    {}
    [\vspace{2pt}\color{border}\hrule height 0.8pt\vspace{6pt}]

\titleformat{\subsection}
    {\color{primary}\fontsize{14}{18}\bfseries\sffamily}
    {\color{accent}\thesubsection}
    {0.8em}
    {}

\titlespacing*{\section}{0pt}{20pt plus 4pt minus 2pt}{8pt plus 2pt minus 2pt}
\titlespacing*{\subsection}{0pt}{16pt plus 3pt minus 2pt}{6pt plus 2pt minus 1pt}

% ========== BASIC HEADERS AND FOOTERS ==========
\usepackage{fancyhdr}
\usepackage{lastpage}
\pagestyle{fancy}
\fancyhf{}
\renewcommand{\headrulewidth}{0.5pt}
\renewcommand{\footrulewidth}{0pt}
\fancyhead[L]{\color{textgray}\small\sffamily Your Document Title Here}
\fancyhead[R]{\color{textgray}\small\sffamily\thepage}
\fancyfoot[C]{\color{textgray}\footnotesize\sffamily Page \thepage{} of \pageref{LastPage}}

% ========== DOCUMENT METADATA ==========
\usepackage{titling}
\pretitle{
    \begin{center}
    \vspace{0.5in}
    \color{primary}\Huge\bfseries\sffamily
}
\posttitle{
    \end{center}
    \vspace{0.3in}
    \begin{center}
    \color{border}\rule{0.6\textwidth}{2pt}
    \end{center}
    \vspace{0.2in}
}
\preauthor{\begin{center}\color{secondary}\large\sffamily}
\postauthor{\end{center}}
\predate{\begin{center}\color{textgray}\large\sffamily}
\postdate{\end{center}\vspace{0.5in}}

\title{Your Document Title Here}
\author{Author Name}
\date{\today}


% ========== DOCUMENT BEGIN ==========
\begin{document}

% ========== COVER PAGE (OPTIONAL) ==========
% Include selected cover if one is specified
\ifthenelse{\equal{\covertype}{none}}{
    % No cover page selected, continue with regular title page
}{
    % Create cover page based on selected cover type
    \clearpage
    \thispagestyle{empty}
    % Handle cover1.tex
    \ifthenelse{\equal{\covertype}{1}}{
        % Load TikZ libraries needed for cover1
        \usetikzlibrary{positioning,shapes.geometric,decorations.pathmorphing,shadows.blur,backgrounds,calc}
        
        % Define placeholder commands if not already defined
        \providecommand{\TitlePlaceholder}{\@title}
        \providecommand{\SummaryTitlePlaceholder}{Executive Summary}
        \providecommand{\SummaryTextPlaceholder}{This document provides comprehensive information on our latest research and development efforts.}
        \providecommand{\PlatformsList}{Web, Mobile, Desktop, Cloud}
        \providecommand{\FocusList}{Performance, Scalability, Security, Usability}
        \providecommand{\AuthorPlaceholder}{\@author}
        \providecommand{\DatePlaceholder}{\@date}
        \providecommand{\ClassificationPlaceholder}{\documentClassification}
        \providecommand{\DistributionPlaceholder}{All Departments}
        \providecommand{\VersionPlaceholder}{\documentVersion}
        \providecommand{\StatusPlaceholder}{\documentStatus}
        
        % Background gradient and shapes
        \begin{tikzpicture}[remember picture,overlay]
            \fill[top color=primary, bottom color=primary!80] (current page.north west) rectangle (current page.south east);
            \fill[accent!20, opacity=0.3] (current page.north east) -- ++(-8cm,0) -- ++(0,-12cm) -- ++(8cm,0) -- cycle;
            \fill[lightaccent!15, opacity=0.4] (current page.south west) -- ++(6cm,0) -- ++(0,10cm) -- ++(-6cm,0) -- cycle;
            \fill[accent!25, opacity=0.5] ([xshift=3cm, yshift=-5cm]current page.north east) circle (4cm);
            \fill[lightaccent!20, opacity=0.6] ([xshift=-4cm, yshift=6cm]current page.south west) circle (3cm);
            \foreach \x in {-10,-8,...,30} {
                \foreach \y in {-20,-18,...,35} {
                    \fill[white, opacity=0.03] (\x,\y) circle (0.5mm);
                }
            }
        \end{tikzpicture}

        % Content overlay
        \begin{tikzpicture}[remember picture,overlay]
            % Title Section
            \node[anchor=north, yshift=-3cm] at (current page.north) {
                \begin{minipage}{0.85\paperwidth}
                    \centering
                    \begin{spacing}{3.5}
                        \textcolor{white}{\fontsize{36}{60}\selectfont\bfseries\sffamily \TitlePlaceholder}
                    \end{spacing}
                \end{minipage}
            };

            % Summary Box
            \node[anchor=center, yshift=1cm] at (current page.center) {
                \begin{minipage}{0.75\paperwidth}
                    \begin{tcolorbox}[
                        colback=secondary, opacityback=0.45,
                        colframe=accent!60, opacityframe=0.45,
                        coltext=white, boxrule=1pt, arc=8pt,
                        left=25pt, right=25pt, top=20pt, bottom=20pt,
                        enhanced, drop shadow, borderline={2pt}{0pt}{accent!60}
                    ]
                    \centering
                    \textcolor{white!95}{\fontsize{14}{18}\selectfont\bfseries\sffamily \SummaryTitlePlaceholder}
                    \vspace{12pt}
                    \begin{spacing}{1.5}
                        \textcolor{white!95}{\fontsize{12}{16}\selectfont\sffamily \SummaryTextPlaceholder}
                    \end{spacing}
                    \vspace{16pt}
                    \begin{minipage}{0.48\textwidth}
                        \textcolor{accent!80}{\fontsize{11}{14}\selectfont\bfseries\sffamily PLATFORMS ANALYZED}
                        \textcolor{white!90}{\fontsize{10}{13}\selectfont\sffamily
                        \PlatformsList
                        }
                    \end{minipage}
                    \hfill
                    \begin{minipage}{0.48\textwidth}
                        \textcolor{accent!80}{\fontsize{11}{14}\selectfont\bfseries\sffamily KEY FOCUS AREAS}
                        \textcolor{white!90}{\fontsize{10}{13}\selectfont\sffamily
                        \FocusList
                        }
                    \end{minipage}
                    \end{tcolorbox}
                \end{minipage}
            };
            
            % Footer Section
            \node[anchor=south, yshift=2.5cm] at (current page.south) {
                \begin{minipage}{0.85\paperwidth}
                    \centering
                    \textcolor{white!90}{\fontsize{16}{20}\selectfont\bfseries\sffamily \AuthorPlaceholder}
                    \vspace{0.4cm}
                    
                    \textcolor{accent!60}{\fontsize{14}{18}\selectfont\sffamily \DatePlaceholder}
                    \vspace{0.8cm}
                    \begin{tcolorbox}[
                        colback=secondary, opacityback=0.45,
                        colframe=accent!60, opacityframe=0.45,
                        coltext=white!90, boxrule=1pt, arc=4pt,
                        left=20pt, right=20pt, top=12pt, bottom=12pt,
                        width=0.7\paperwidth, center
                    ]
                    \centering
                    \fontsize{10}{12}\selectfont\sffamily
                    \textbf{DOCUMENT CLASSIFICATION:} \ClassificationPlaceholder \quad
                    \textbf{DISTRIBUTION:} \DistributionPlaceholder
                    \\
                    \textbf{VERSION:} \VersionPlaceholder \quad \textbf{STATUS:} \StatusPlaceholder
                    \end{tcolorbox}
                \end{minipage}
            };
            
            % Watermark
            \node[anchor=center, opacity=0.05] at (current page.center) {
                \textcolor{white}{\fontsize{120}{140}\selectfont\bfseries\sffamily \TitlePlaceholder}
            };
        \end{tikzpicture}
    }{}
    
    % Handle cover2.tex
    \ifthenelse{\equal{\covertype}{2}}{
        % Additional TikZ libraries needed for cover2
        \usetikzlibrary{positioning,shapes.geometric,decorations.pathmorphing}
        
        % Check if we have FP package loaded, if not load it
        \makeatletter
        \@ifundefined{fpeval}{
            \usepackage{fp}
        }{}
        \makeatother
        
        % Define cover2 specific commands
        \newcommand{\coverrule}{%
            \begin{center}
            \begin{tikzpicture}
                \draw[accent, line width=2pt] (0,0) -- (4,0);
                \draw[lightaccent, line width=1pt] (4.5,0) -- (6,0);
                \draw[border, line width=0.5pt] (6.5,0) -- (7.5,0);
            \end{tikzpicture}
            \end{center}
        }
        
        \newcommand{\covertitle}[1]{%
            \begin{center}
            \color{primary}
            \fontsize{32}{38}\selectfont\bfseries\sffamily
            #1
            \end{center}
        }
        
        \newcommand{\coversubtitle}[1]{%
            \begin{center}
            \color{secondary}
            \fontsize{18}{22}\selectfont\sffamily
            #1
            \end{center}
        }
        
        \newcommand{\coverauthor}[1]{%
            \begin{center}
            \color{textgray}
            \fontsize{14}{18}\selectfont\sffamily
            #1
            \end{center}
        }
        
        \newtcolorbox{coverinfobox}{
            colback=background,
            colframe=accent,
            coltext=primary,
            boxrule=1.5pt,
            arc=6pt,
            left=20pt,
            right=20pt,
            top=16pt,
            bottom=16pt,
            enhanced,
            drop shadow,
            width=0.85\textwidth,
            center
        }
        
        \newtcolorbox{classificationbox}{
            colback=lightgray,
            colframe=border,
            coltext=textgray,
            boxrule=0.8pt,
            arc=3pt,
            left=12pt,
            right=12pt,
            top=8pt,
            bottom=8pt,
            width=0.6\textwidth,
            center
        }
        
        % Define placeholder commands if not already defined
        \providecommand{\TitlePlaceholder}{\@title}
        \providecommand{\SubtitlePlaceholder}{Professional documentation system}
        \providecommand{\SummaryTitlePlaceholder}{Executive Summary}
        \providecommand{\SummaryTextPlaceholder}{This document provides comprehensive information on our latest research and development efforts.}
        \providecommand{\FocusItemOne}{Technical Excellence}
        \providecommand{\FocusItemTwo}{User Experience}
        \providecommand{\FocusItemThree}{Scalability}
        \providecommand{\FocusItemFour}{Market Advantage}
        \providecommand{\AuthorPlaceholder}{\@author}
        \providecommand{\DatePlaceholder}{\@date}
        \providecommand{\ClassificationPlaceholder}{\documentClassification}
        \providecommand{\DistributionPlaceholder}{All Departments}
        \providecommand{\VersionPlaceholder}{\documentVersion}
        \providecommand{\ClassificationLevelPlaceholder}{Confidential}
        
        % Cover 2 content
        \vspace*{0.5in}
        
        % Decorative circle element
        \begin{center}
        \begin{tikzpicture}
            \fill[accent!15] (0,0) circle (1.2);
            \fill[accent!30] (0,0) circle (0.8);
            \fill[accent] (0,0) circle (0.4);
            \fill[lightaccent!20] (-3,0) circle (0.3);
            \fill[lightaccent!20] (3,0) circle (0.3);
            \draw[accent!40, line width=1pt] (-2.7,0) -- (-1.2,0);
            \draw[accent!40, line width=1pt] (1.2,0) -- (2.7,0);
        \end{tikzpicture}
        \end{center}
        
        \vspace{0.4in}
        
        % Title and subtitle
        \covertitle{\TitlePlaceholder}
        
        \vspace{0.2in}
        \coverrule
        \vspace{0.3in}
        
        \coversubtitle{\SubtitlePlaceholder}
        
        \vspace{0.6in}
        
        \begin{coverinfobox}
        \textbf{\large \SummaryTitlePlaceholder}
        
        \vspace{8pt}
        
        \SummaryTextPlaceholder
        
        \vspace{8pt}
        
        \textbf{Key Focus Areas:}
        \begin{itemize}
            \item \FocusItemOne
            \item \FocusItemTwo
            \item \FocusItemThree
            \item \FocusItemFour
        \end{itemize}
        \end{coverinfobox}
        
        \vspace{0.8in}
        
        \coverauthor{\AuthorPlaceholder}
        
        \vspace{0.2in}
        
        \begin{center}
        \color{accent}
        \fontsize{16}{20}\selectfont\bfseries\sffamily
        \DatePlaceholder
        \end{center}
        
        \vfill
        
        \begin{center}
        \begin{tikzpicture}
            \foreach \i in {0,0.2,...,6} {
                \draw[accent!\fpeval{100-\i*10}, line width=\fpeval{2-\i*0.2}pt] 
                (\i*0.5,0) -- (\i*0.5+0.3,0);
            }
        \end{tikzpicture}
        \end{center}
        
        \vspace{0.3in}
        
        \begin{classificationbox}
        \centering
        \footnotesize\sffamily
        \textbf{Document Classification:} \ClassificationPlaceholder\\
        \textbf{Distribution:} \DistributionPlaceholder\\
        \textbf{Version:} \VersionPlaceholder \quad 
        \textbf{Classification:} \ClassificationLevelPlaceholder
        \end{classificationbox}
    }{}
    
    \newpage
}

% ========== TITLE PAGE ==========
% Instructions: This creates the title page automatically
\maketitle
\thispagestyle{empty}

% ========== EXECUTIVE SUMMARY BOX ==========
% Instructions: Replace with your document summary or overview
\begin{infobox}
\textbf{From \brandname{\companyName}}: This document demonstrates the advanced typesetting capabilities of our \productname{Lagen} documentation system. It showcases professional formatting features including multi-page tables, enhanced column types, advanced mathematical typesetting, and sophisticated branding controls.
\end{infobox}

\vfill
% ========== DOCUMENT CLASSIFICATION ==========
% Instructions: Replace with your document's classification information
\begin{center}
\color{textgray}\footnotesize\sffamily
\textbf{Document Classification:} \documentClassification\\
\textbf{Distribution:} All Departments\\
\textbf{Version:} \documentVersion \quad \textbf{Pages:} \pageref{LastPage}
\end{center}

\newpage

% ========== TABLE OF CONTENTS ==========
% Instructions: This generates automatically based on your sections
\tableofcontents
\newpage

% ========== MAIN CONTENT SECTIONS ==========

\section{Content Demonstration}
\label{sec:demonstration}
% Instructions: Replace with your first major section title

\lettrine[lines=3, loversize=0.1]{T}{his} demonstrates a drop cap at the beginning of a section using the lettrine package. It adds an elegant typographic touch to your document. The rest of this paragraph shows how the text wraps around the large initial letter.

This document now features \nameref{sec:tables} spanning multiple pages, proper formatting of units like \SI{42}{\kilo\watt\hour} or \SI{1.25}{\USD}, and advanced mathematical typesetting with $\displaystyle\int_{0}^{\infty} e^{-x^2} dx = \frac{\sqrt{\pi}}{2}$.

\subsection{Advanced Typography Features}

\needspace{4\baselineskip}
Our branding is consistent with \brandname{\companyName} and \productname{Lagen Documentation System} appearing throughout the document. 

Here we demonstrate quotation handling with csquotes: "This is properly formatted as a quotation with smart quotes." Notice how the quotes are properly curled rather than using straight quotes.

\begin{displayquote}
This is a block quotation that shows proper indentation and formatting. It's useful for longer quotes that deserve special emphasis or visual distinction from the main text.
\end{displayquote}

\subsection{Cross-References and Citations}

This section demonstrates intelligent cross-referencing. You can refer to Section~\ref{sec:boxes} for box demonstrations, or Section~\nameref{sec:tables} for table examples.

\needspace{4\baselineskip}
For bibliography management, you can cite sources like \texttt{\textbackslash cite\{example\}} and have them automatically formatted based on your chosen citation style.\footnote{Requires a .bib file and proper setup with biblatex and biber.}

\subsection{Subsection Example}
% Instructions: Use subsections to organize content within main sections

This demonstrates how to organize content using subsections.

\begin{itemize}
    \item This is a first-level bullet point
    \item This is another bullet point with nested items:
    \begin{itemize}
        \item This is a second-level bullet point
        \item This is another second-level bullet point
    \end{itemize}
\end{itemize}

\subsubsection{Nested Subsection}
% Instructions: Use subsubsections for further detail within subsections

This shows a third-level heading for detailed organization.

\section{Table Examples}
\label{sec:tables}
% Instructions: This section demonstrates modern table formatting

This section shows how to create professional-looking tables with color and styling, now with multi-page support and enhanced column types.

\subsection{Enhanced Table Examples}

The following demonstrates a table with custom column types and multirow capabilities:

\begin{table}[htbp]
\caption{Advanced Table Formatting}
\centering
\rowcolors{2}{rowalt}{white}
\begin{tabular}{L{2.5cm}C{4cm}R{2cm}N}
\rowcolor{headerbg}
\textcolor{headertext}{\textbf{Category}} &
\textcolor{headertext}{\textbf{Description}} &
\textcolor{headertext}{\textbf{Value}} &
\textcolor{headertext}{\textbf{Progress}} \\
\toprule
\multirow{2}{*}{Development} & Initial framework design & \SI{15000}{\USD} & 95\% \\
 & Component integration & \SI{8500}{\USD} & 45\% \\
\midrule
\multirow{3}{*}{Testing} & Unit testing framework & \SI{4200}{\USD} & 100\% \\
 & Integration testing & \SI{6300}{\USD} & 60\% \\
 & Performance testing & \SI{5500}{\USD} & 30\% \\
\midrule
Deployment & Cloud infrastructure & \SI{12000}{\USD} & 15\% \\
\bottomrule
\end{tabular}
\end{table}

\subsection{Multi-page Table Example}

Here's an example of a table that can span multiple pages:

\begin{longtable}{p{2cm}p{4cm}p{3cm}p{4cm}}
\caption{Multi-page Project Timeline and Milestones} \\
\rowcolor{headerbg}
\textcolor{headertext}{\textbf{Date}} &
\textcolor{headertext}{\textbf{Milestone}} &
\textcolor{headertext}{\textbf{Responsible}} &
\textcolor{headertext}{\textbf{Deliverables}} \\
\toprule
\endfirsthead

\multicolumn{4}{c}{\tablename\ \thetable{} -- continued from previous page} \\
\rowcolor{headerbg}
\textcolor{headertext}{\textbf{Date}} &
\textcolor{headertext}{\textbf{Milestone}} &
\textcolor{headertext}{\textbf{Responsible}} &
\textcolor{headertext}{\textbf{Deliverables}} \\
\toprule
\endhead

\midrule \multicolumn{4}{r}{Continued on next page} \\
\endfoot

\bottomrule
\endlastfoot

2023-01-15 & Project Initiation & Project Manager & Project charter, stakeholder register \\
\rowcolor{rowalt}
2023-02-01 & Requirements Gathering & Business Analyst & Requirements document, user stories \\
2023-02-15 & System Architecture & Solution Architect & Architecture document, technology stack \\
\rowcolor{rowalt}
2023-03-01 & Design Phase Completion & Design Team & UI/UX designs, wireframes \\
2023-03-15 & Development Sprint 1 & Development Team & Core functionality implementation \\
\rowcolor{rowalt}
2023-04-01 & Development Sprint 2 & Development Team & Feature set 1 completion \\
2023-04-15 & Development Sprint 3 & Development Team & Feature set 2 completion \\
\rowcolor{rowalt}
2023-05-01 & Alpha Testing & QA Team & Test reports, bug tracking \\
2023-05-15 & Beta Release & Release Manager & Beta deployment, user feedback \\
\rowcolor{rowalt}
2023-06-01 & User Acceptance Testing & QA Team & UAT report, sign-off documents \\
2023-06-15 & Performance Optimization & Performance Engineer & Optimization report, benchmarks \\
\rowcolor{rowalt}
2023-07-01 & Security Audit & Security Team & Security assessment, vulnerability report \\
2023-07-15 & Documentation & Technical Writer & User manual, API documentation \\
\rowcolor{rowalt}
2023-08-01 & Final QA & QA Team & Final test report, release readiness \\
2023-08-15 & Production Deployment & DevOps Team & Deployment checklist, monitoring setup \\
\rowcolor{rowalt}
2023-09-01 & Post-Launch Support & Support Team & Support plan, escalation procedures \\
2023-09-15 & Project Review & Project Manager & Lessons learned, project closure report \\
\end{longtable}

\subsection{Data Analysis Table}

\begin{tcolorbox}[colback=cardbg, colframe=frameborder, coltitle=white,
    fonttitle=\bfseries, title=📊 Data Analysis Results,
    enhanced, boxrule=0.5pt, arc=4pt, left=10pt, right=10pt, top=8pt, bottom=8pt,
    attach boxed title to top left={xshift=10pt, yshift=-2mm},
    boxed title style={colback=accent, arc=2pt}
]

\rowcolors{2}{rowalt}{white}

\begin{tabularx}{\linewidth}{>{\raggedright\arraybackslash\bfseries}X >{\raggedright\arraybackslash}X c c}
    \rowcolor{headerbg}
    \textcolor{headertext}{Category} &
    \textcolor{headertext}{Description} &
    \textcolor{headertext}{Value} &
    \textcolor{headertext}{Status} \\
    \toprule
    Project A & Implementation phase & 95\% & Complete  \\
    Project B & Analysis phase & 45\% & In progress \\
    Project C & Planning phase & 20\% & Started \\
    Project D & Testing phase & 80\% & Advanced \\
    Project E & Deployment phase & 10\% & Initial \\
    \bottomrule
\end{tabularx}

\end{tcolorbox}

\section{Mathematical Examples}
\label{sec:math}

\subsection{Equation Examples}

The enhanced math support allows for beautifully typeset equations:

\begin{align}
E &= mc^2 \\
F &= G\frac{m_1 m_2}{r^2} \\
\nabla \times \vec{E} &= -\frac{\partial \vec{B}}{\partial t}
\end{align}

\subsection{Units and Numbers}

With the siunitx package, we can consistently format units:

\begin{itemize}
    \item The battery capacity is \SI{5000}{\milli\ampere\hour}
    \item The system requires \SI{240}{\volt} power
    \item The project will take approximately \SI{3.5}{\year} to complete
    \item The budget is \SI{2500000}{\USD} for the first phase
    \item Normal room temperature is \SI{20}{\celsius}
    \item The server processes \SI{5.2}{\giga\byte} of data per minute
\end{itemize}

\section{Pseudocode and Algorithms}
\label{sec:algorithms}

\subsection{Basic Algorithm Example}

Here is a simple algorithm example using the algorithm2e package:

\begin{algorithm}[H]
	\SetAlgoLined
	\KwData{Input data $X$, parameters $\alpha$, $\beta$}
	\KwResult{Optimized output $Y$}
	Initialize $Y \gets \emptyset$\;
	\For{each element $x \in X$}{
		$y \gets \text{Process}(x, \alpha)$\;
		\If{$y > \beta$}{
			$Y \gets Y \cup \{y\}$\;
		}
		\Else{
			$y \gets \text{AlternativeProcess}(x)$\;
			$Y \gets Y \cup \{y\}$\;
		}
	}
	\Return{$Y$}\;
	\caption{Basic Optimization Algorithm}
\end{algorithm}

\section{Advanced Box Types}
\label{sec:boxes}
% Instructions: Examples of how to use the three types of callout boxes

Here are examples of different types of callout boxes:

\begin{successbox}
\textbf{Success Box:} Use this for highlighting positive outcomes or important accomplishments.
\end{successbox}

\begin{alertbox}
\textbf{Alert Box:} Use this for important cautions or critical information.
\end{alertbox}

\begin{infobox}[title=Information Box with Title]
This demonstrates an information box with a title. Use for supplementary information.
\end{infobox}

\subsection{Icon-Enhanced Boxes}

The awesomebox package provides boxes with built-in icons:

\awesomebox[violet]{2pt}{\faInfo}{violet}{This is a note with an icon. It's useful for additional information that deserves attention.}

\awesomebox[teal]{2pt}{\faLightbulb}{teal}{This is a tip with an icon. It provides helpful suggestions or shortcuts.}

\awesomebox[orange]{2pt}{\faExclamationTriangle}{orange}{This is a warning with an icon. It alerts readers to potential issues or mistakes.}

\awesomebox[red]{2pt}{\faExclamationCircle}{red}{This is a caution with an icon. It emphasizes important concerns.}

\awesomebox[blue]{2pt}{\faInfoCircle}{blue}{This is an important box with an icon. It highlights critical information.}

\subsection{Logo-Enhanced Boxes}

The bclogo package offers boxes with logos:

\bclogo[logo=\bcattention]{This box has an attention logo to alert readers to something important.}

\bclogo[logo=\bcquestion]{This box has a question logo for FAQs or exploration questions.}

\bclogo[logo=\bccrayon]{This box has a crayon logo that might be used for creative exercises or notes.}

\section{Code Listings and Algorithms}

\subsection{Code Snippets with Syntax Highlighting}

Here's an example of code with syntax highlighting using the listings package:

\begin{lstlisting}[language=Python, caption=Sample Python Code]
# This is a comment
def hello_world():
    """Docstring for function"""
    print("Hello, world!")
    return True

# Call the function
result = hello_world()
if result:
    print("Function executed successfully")
\end{lstlisting}

% Uncomment this if you have minted package configured
% \begin{minted}[frame=single, framesep=8pt]{javascript}
% // This is a comment
% function helloWorld() {
%   // Print a message
%   console.log("Hello, world!");
%   return true;
% }
% 
% // Call the function
% const result = helloWorld();
% if (result) {
%   console.log("Function executed successfully");
% }
% \end{minted}

\subsection{Algorithm Description}

Here's an example of an algorithm description using the algorithm2e package:

\begin{algorithm}[H]
\caption{Sample Algorithm}
\SetAlgoLined
\KwData{Input data $X$}
\KwResult{Output result $Y$}
Initialize parameters\;
\For{each element in $X$}{
    Process the element\;
    \If{condition is met}{
        Perform special operation\;
    }
    Update intermediate results\;
}
Calculate final result $Y$\;
\Return{$Y$}\;
\end{algorithm}

\section{List Examples}
% Instructions: Examples of different list types

\subsection{Bullet Point Lists}
% Instructions: Standard bullet points for unordered information

\begin{itemize}
    \item \textbf{First point:} Explanation with emphasized beginning.
    
    \item \textbf{Second point:} Another example with supporting detail.
\end{itemize}

\subsection{Numbered Lists}
% Instructions: Use enumerate environment for ordered/sequential information

\begin{enumerate}
    \item \textbf{First step:} Explanation of the first step in a process.
    
    \item \textbf{Second step:} Details about the second step or item.
\end{enumerate}

\section{Analysis Example}

This section shows how to structure content with multiple levels.

\begin{infobox}[title=Key Points]
\textbf{Point 1:} First key point\\
\textbf{Point 2:} Second key point\\
\textbf{Point 3:} Third key point
\end{infobox}

\section{References}
% Instructions: Section for citations and resources

% Example bibliography entries for demonstration
\begin{filecontents*}{example.bib}
@article{example2023,
  author = {Smith, John},
  title = {Example Title},
  year = {2023},
  journal = {Journal of Examples},
  volume = {10},
  number = {2},
  pages = {100--120}
}
\end{filecontents*}
% Note: In a real document, you would place biblatex entries in a separate .bib file

\textbf{References:}
\begin{itemize}
    \item \href{https://example.com/source1}{Reference 1 - Brief Description}
    \item \href{https://example.com/source2}{Reference 2 - Brief Description}
\end{itemize}

% Uncomment to print bibliography when you have a .bib file
% \printbibliography[heading=bibintoc, title={Bibliography}]

% Instructions: This label is used for page counting in the footer
\label{LastPage}
\end{document}
